\chapter{函数极限：习题解答}

\section*{A组　定义、量词与基本性质}
\addcontentsline{toc}{section}{A组　定义、量词与基本性质}

\begin{solution}{14.1}
\(\{1/n:n\in\Npos\}\) 的唯一聚点是 \(0\)。任意 \(\delta>0\) 可取 \(n>1/\delta\)，使 \(0<1/n<\delta\)。若 \(a\ne0\)，在 \(a<0\) 时可取不与正半轴相交的小邻域；在 \(a>0\) 时，序列 \(1/n\) 在 \(a\) 附近只有有限项，因而可缩小邻域避开它们。

\(\Q\cap[0,1]\) 的聚点集为 \([0,1]\)：有理数在每个非退化实区间中稠密，端点用单侧点即可；区间外的点与 \([0,1]\) 距离为正。\(\Z\) 没有聚点，因为对每个 \(m\in\Z\) 可取半径 \(1/2\) 排除其余整数，非整数点也可取小邻域至多含一个整数。
\end{solution}

\begin{solution}{14.2}
若 \(a\) 是孤立点，存在 \(\delta_0>0\) 使
\[
 D\cap U^\circ(a,\delta_0)=\varnothing.
\]
于是对任意 \(L\) 和任意 \(\varepsilon>0\)，选 \(\delta\le\delta_0\) 后，定义中的蕴含式因前件从不成立而为真。这样每个 \(L\) 都会成为“极限”，与唯一性冲突。聚点条件排除了这种空泛成立。
\end{solution}

\begin{solution}{14.3}
设 \(g(x)=f(x)\) 对所有 \(x\in D\setminus\{a\}\) 成立，而 \(f(a),g(a)\) 可以不同。极限定义只检验 \(0<\abs{x-a}<\delta\) 的点，因此在这些点上
\[
 \abs{g(x)-L}=\abs{f(x)-L}.
\]
同一个 \(\delta(\varepsilon)\) 同时适用于两函数，故二者在 \(a\) 处有相同极限。
\end{solution}

\begin{solution}{14.4}
严格形式为
\[
 \forall\varepsilon>0\ \exists\delta>0\ \forall x\in D,\quad
 0<\abs{x-a}<\delta\Rightarrow\abs{f(x)-L}<\varepsilon.
\]
否定时依次交换全称与存在量词，并否定蕴含式，得到
\[
 \exists\varepsilon_0>0\ \forall\delta>0\ \exists x\in D,\quad
 0<\abs{x-a}<\delta,\quad\abs{f(x)-L}\ge\varepsilon_0.
\]
\end{solution}

\begin{solution}{14.5}
给定 \(\varepsilon>0\)，取 \(\delta=\varepsilon/2\)。若 \(0<\abs{x-3}<\delta\)，则
\[
 \abs{(2x-5)-1}=2\abs{x-3}<2\delta=\varepsilon.
\]
故极限为 \(1\)。
\end{solution}

\begin{solution}{14.6}
分解
\[
 \abs{x^3-a^3}=\abs{x-a}\abs{x^2+ax+a^2}.
\]
先要求 \(\abs{x-a}<1\)，则 \(\abs x\le\abs a+1\)，所以
\[
 \abs{x^2+ax+a^2}
 \le(\abs a+1)^2+\abs a(\abs a+1)+a^2=:C_a.
\]
这里 \(C_a>0\)。给定 \(\varepsilon>0\)，取
\[
 \delta=\min\{1,\varepsilon/C_a\}.
\]
于是目标误差小于 \(\varepsilon\)。
\end{solution}

\begin{solution}{14.7}
反三角不等式给出
\[
 \bigl|\abs{f(x)}-\abs L\bigr|\le\abs{f(x)-L}.
\]
对给定 \(\varepsilon>0\)，沿用 \(f(x)\to L\) 对应的半径，右端小于 \(\varepsilon\)，故 \(\abs{f(x)}\to\abs L\)。
\end{solution}

\begin{solution}{14.8}
在极限定义中取 \(\varepsilon=L/2\)。存在 \(\delta>0\) 使删心邻域内
\[
 \abs{f(x)-L}<L/2.
\]
因而 \(f(x)>L-L/2=L/2\)。条件 \(L>0\) 保证所选误差为正并给出正下界。
\end{solution}

\begin{solution}{14.9}
设在 \(a\) 的某删心邻域内 \(f(x)\le h(x)\le g(x)\)，且 \(f(x),g(x)\to L\)。给定 \(\varepsilon>0\)，取共同半径，使
\[
 L-\varepsilon<f(x),\qquad g(x)<L+\varepsilon
\]
同时成立，再与次序合并：
\[
 L-\varepsilon<h(x)<L+\varepsilon.
\]
故 \(h(x)\to L\)。
\end{solution}

\begin{solution}{14.10}
若题设中的点列 \((x_n)\subset D\setminus\{a\}\) 确实趋于 \(a\)，并且 \(f,g\) 的函数极限都存在，则保序结论仍成立。记极限为 \(L,M\)。若 \(L>M\)，取 \(\varepsilon=(L-M)/3\)。由两个函数极限和 \(x_n\to a\)，充分大的 \(n\) 满足
\[
 f(x_n)>L-\varepsilon>M+\varepsilon>g(x_n),
\]
与 \(f(x_n)\le g(x_n)\) 矛盾。

这条较弱假设只足以比较两个已经存在的极限，并不能保证极限存在。例如 \(f(x)=\sin(1/x)\)、\(g(x)=1\) 在 \(x_n=1/(2\pi n)\) 上满足 \(f(x_n)\le g(x_n)\)，但 \(f\) 在 \(0\) 处无极限。若点列不趋于 \(a\)，则连极限次序也无法推出。
\end{solution}

\begin{solution}{14.11}
有
\[
 \abs{x^2+3x-10}
 =\abs{x-2}\abs{x+5}.
\]
若 \(\abs{x-2}<1\)，则 \(\abs{x+5}\le\abs{x-2}+7<8\)。给定 \(\varepsilon>0\)，取
\[
 \delta=\min\{1,\varepsilon/8\},
\]
便有目标误差小于 \(8\delta\le\varepsilon\)。
\end{solution}

\begin{solution}{14.12}
对 \(x\ge0\)，
\[
 \abs{\sqrt x-2}
 =\frac{\abs{x-4}}{\sqrt x+2}
 \le\frac12\abs{x-4}.
\]
给定 \(\varepsilon>0\)，取 \(\delta=2\varepsilon\)，即得 \(\abs{\sqrt x-2}<\varepsilon\)。
\end{solution}

\begin{solution}{14.13}
设在某删心邻域内 \(\abs{g(x)}\le M\)。若 \(M=0\)，乘积恒为零；若 \(M>0\)，给定 \(\varepsilon>0\)，由 \(f(x)\to0\) 取半径使
\[
 \abs{f(x)}<\varepsilon/M.
\]
再与有界性半径取最小值，便有 \(\abs{f(x)g(x)}<\varepsilon\)。
\end{solution}

\begin{solution}{14.14}
均在 \(x\to0\) 时考虑，取 \(f(x)=x\)。令
\[
 g_1(x)=\frac1{\sqrt{\abs x}},\qquad
 g_2(x)=\frac1x,\qquad
 g_3(x)=\frac{\sin(1/x)}x.
\]
三者在每个删心邻域内无界，而
\[
 f g_1=\operatorname{sgn}(x)\sqrt{\abs x}\to0,\quad
 f g_2=1,\quad
 f g_3=\sin(1/x)
\]
分别趋于零、趋于 \(1\)、无极限。最后一项沿 \(x_n=1/(\pi/2+2\pi n)\) 与 \(y_n=1/(3\pi/2+2\pi n)\) 分别取值 \(1,-1\)。
\end{solution}

\begin{solution}{14.15}
设 \(g(x)\to M\ne0\)。取误差 \(\abs M/2\)，可得某删心邻域内
\[
 \abs{g(x)}>\abs M/2,
\]
故倒数有定义。再由
\[
 \abs{\frac1{g(x)}-\frac1M}
 =\frac{\abs{g(x)-M}}{\abs{g(x)}\abs M}
 \le\frac2{\abs M^2}\abs{g(x)-M}
\]
知倒数趋于 \(1/M\)。非零条件既使 \(1/M\) 有定义，又产生分母的正下界。
\end{solution}

\begin{solution}{14.16}
若 \(q(a)\ne0\)，商法则直接给出有限极限，例如
\[
 \frac{x+1}{x+2}\longrightarrow\frac23\quad(x\to1).
\]
若 \(q(a)=0\) 且 \(p\) 不恒为零，把两多项式写成
\[
 p(x)=(x-a)^mP(x),\qquad q(x)=(x-a)^nQ(x),
\]
其中 \(P(a)Q(a)\ne0\)，并允许 \(m=0\)。若 \(p\equiv0\)，商在有定义处恒为零。对非零 \(p\)，当 \(m\ge n\) 时，公共零因子可把分母在 \(a\) 的零点完全消去；若 \(m=n\)，极限为 \(P(a)/Q(a)\)，若 \(m>n\)，极限为零。例如
\[
 \frac{x^2-1}{x-1}=x+1\to2\quad(x\to1).
\]
当 \(m<n\) 时，消去公共因子后分母仍含 \((x-a)^{n-m}\)。此时商的绝对值趋于无穷；剩余次数为偶数时两侧符号相同，为奇数时两侧符号相反，具体符号还由 \(P(a)/Q(a)\) 决定。例如
\[
 \frac1{x^2}\to+\infty,\qquad \frac1x
\]
在 \(x\to0\) 时左右符号相反。分类所需的是分子、分母在 \(a\) 的零点阶数，而不是“可约”这一含混表述。
\end{solution}

\begin{solution}{14.17}
令 \(b=0\)、\(g(x)\equiv0\)。定义
\[
 f(y)=
 \begin{cases}
 y,&y\ne0,\\
 1,&y=0.
 \end{cases}
\]
则 \(g(x)\to0\)、\(\lim_{y\to0}f(y)=0\)，但 \(f(g(x))\equiv1\)。内层函数持续取中心值，而外层删心极限不控制该点值。
\end{solution}

\begin{solution}{14.18}
设 \(f\) 在 \(b\) 连续。给定 \(\varepsilon>0\)，存在 \(\eta>0\) 使
\[
 y\in E,\quad\abs{y-b}<\eta
 \Longrightarrow\abs{f(y)-f(b)}<\varepsilon.
\]
这里允许 \(y=b\)，因为此时误差为零。由 \(g(x)\to b\) 取 \(\delta\)，使 \(0<\abs{x-a}<\delta\) 时 \(\abs{g(x)-b}<\eta\)，直接代入即可。
\end{solution}

\begin{solution}{14.19}
任意 \(\delta>0\) 可取 \(n>1/\delta\)，使 \(1/n\in D\) 且 \(0<1/n<\delta\)，故 \(0\) 是聚点。若极限为 \(L\)，偶数指标与奇数指标给出的值分别为 \(1,-1\)。取
\[
 \varepsilon_0=\frac12.
\]
若 \(L\ge0\)，在任意邻域选足够大的奇数 \(n\)，则 \(\abs{-1-L}\ge1\)；若 \(L<0\)，选足够大的偶数 \(n\)，则 \(\abs{1-L}>1\)。所有候选 \(L\) 均被排除。
\end{solution}

\begin{solution}{14.20}
设 \(\lim_{x\to a}f(x)=L\)，并给定 \(L\) 的任意邻域半径 \(\varepsilon>0\)。按定义存在 \(\delta>0\)，使所有满足
\[
 x\in D,\qquad0<\abs{x-a}<\delta
\]
的点都有 \(f(x)\in(L-\varepsilon,L+\varepsilon)\)。任何趋于 \(a\) 且最终不等于 \(a\) 的定义域点族，按“趋于”的含义最终进入上述 \(\delta\)-邻域，故对应函数值最终进入给定的 \(L\)-邻域。这里只是直接嵌套两个定义，没有使用逆向的序列刻画。
\end{solution}

\begin{solution}{14.21}
结论一般不成立。令 \(L=1\)，\(f(x)=1\)（\(x\in\Q\)）而 \(f(x)=-1\)（\(x\notin\Q\)），则 \(f^2(x)\equiv1\)，但 \(f\) 在任一点都没有极限。

一个充分条件是 \(f\) 在某删心邻域内保号。若最终 \(f(x)\ge0\)，则
\[
 \abs{f(x)-\abs L}
 =\frac{\abs{f^2(x)-L^2}}{f(x)+\abs L}
\]
在 \(L\ne0\) 时趋于零；若 \(L=0\)，由 \(\abs{f(x)}=\sqrt{f^2(x)}\to0\) 得结论。最终非正的情形类似，极限为 \(-\abs L\)。
\end{solution}

\begin{solution}{14.22}
恒等式
\[
 \max\{u,v\}=\frac{u+v+\abs{u-v}}2,\qquad
 \min\{u,v\}=\frac{u+v-\abs{u-v}}2
\]
把问题化为和、差、数乘以及绝对值的极限运算。若 \(f\to L\)、\(g\to M\)，便有
\[
 \max\{f,g\}\to\max\{L,M\},\qquad
 \min\{f,g\}\to\min\{L,M\}.
\]
\end{solution}

\begin{solution}{14.23}
可按以下四个模块整理。
\begin{enumerate}
 \item 因式分解：把误差写成 \(\abs{x-a}\) 与局部可控因子的乘积；平方、立方和多项式是典型对象。
 \item 有理化：用共轭式把根式差化为自变量差；中心值为零时需另行处理分母下界。
 \item 局部有界因子：由存在有限极限推出局部有界，再控制乘积误差；缺少有界性时，零因子乘无界因子可以产生多种行为。
 \item 分母保离零：非零极限给出统一正下界，从而允许倒数运算；分母极限为零时商法则失效。
\end{enumerate}
每个模块的证明都应先注明固定局部限制承担的作用，再按给定 \(\varepsilon\) 选择最终半径。
\end{solution}

\begin{solution}{14.24}
极限始终相对于定义域。以公式
\[
 f(x)=\frac{\abs x}{x}\qquad(x\ne0)
\]
为例：在 \(D_+=(0,\infty)\) 上令 \(x\to0\)，函数恒为 \(1\)，极限为 \(1\)；在 \(D_-=(-\infty,0)\) 上极限为 \(-1\)；在 \(D=\R\setminus\{0\}\) 上双侧极限不存在。公式相同，允许逼近中心的方向不同，量词所覆盖的点集随之改变。
\end{solution}

\section*{B组　序列刻画与 Cauchy 准则}
\addcontentsline{toc}{section}{B组　序列刻画与 Cauchy 准则}

\begin{solution}{15.1}
若 \(a\) 是聚点，对每个 \(n\) 从
\[
 D\cap\{x:0<\abs{x-a}<1/n\}
\]
中选一点 \(x_n\)，即得 \(x_n\ne a\) 且 \(x_n\to a\)。反之，若有此点列，则给定 \(\delta>0\)，由收敛性可取充分大 \(n\) 使 \(0<\abs{x_n-a}<\delta\)，故每个删心邻域都含定义域点。
\end{solution}

\begin{solution}{15.2}
删心极限不控制 \(f(a)\)。若允许点列恒取 \(x_n=a\)，Heine 条件会强迫 \(f(a)=L\)，从而把连续性条件误加到极限上。例如 \(f(x)=0\)（\(x\ne0\)）、\(f(0)=1\) 的删心极限为 \(0\)，但常值点列 \(x_n=0\) 的像恒为 \(1\)。
\end{solution}

\begin{solution}{15.3}
由聚点性选 \(x_n\in D\setminus\{a\}\) 且 \(x_n\to a\)。若函数极限同时为 \(L,M\)，Heine 定理使同一数列 \(f(x_n)\) 同时趋于 \(L,M\)。数列极限唯一性给出 \(L=M\)。
\end{solution}

\begin{solution}{15.4}
否定形式给出某个 \(\varepsilon_0>0\)，使每个 \(\delta>0\) 内都有坏点。对第 \(n\) 步代入 \(\delta=1/n\)，选
\[
 x_n\in D,\quad0<\abs{x_n-a}<1/n,\quad
 \abs{f(x_n)-L}\ge\varepsilon_0.
\]
第一组不等式推出 \(x_n\to a\)，最后一组不等式表明像数列不可能趋于 \(L\)。
\end{solution}

\begin{solution}{15.5}
若 \(x_n\to a\)，则任意子列 \(x_{n_k}\to a\)，且仍有 \(x_{n_k}\in D\setminus\{a\}\)。对全部逼近点列成立的假设可直接用于这条子列，故 \(f(x_{n_k})\to L\)。也可先由假设得 \(f(x_n)\to L\)，再用数列极限的子列继承性。
\end{solution}

\begin{solution}{15.6}
取 \(x_n=1/n\)、\(y_n=-1/n\)。两列都趋于 \(0\)，但
\[
 \frac{\abs{x_n}}{x_n}=1,\qquad
 \frac{\abs{y_n}}{y_n}=-1.
\]
像数列极限不同，故双侧极限不存在。
\end{solution}

\begin{solution}{15.7}
取
\[
 x_n=\frac1{2\pi n},\qquad y_n=\frac1{(2n+1)\pi}.
\]
则 \(x_n,y_n\to0\)，而
\[
 \cos(1/x_n)=1,\qquad \cos(1/y_n)=-1.
\]
由双路径判别，极限不存在。
\end{solution}

\begin{solution}{15.8}
若 \(x_n,y_n\to a\)，Heine 定理给出
\[
 f(x_n)\to L,\qquad f(y_n)\to L.
\]
数列极限的差法则于是给出 \(f(x_n)-f(y_n)\to0\)。
\end{solution}

\begin{solution}{15.9}
必要性只使用 \(x_n\in D\) 与 \(x_n\to a\)，不要求点列之间填满区间。充分性中若极限失败，每个删心邻域与 \(D\) 的交都含坏点，按半径 \(1/n\) 选取即可。因此证明仅依赖聚点性，适用于任意子集 \(D\subset\R\)。
\end{solution}

\begin{solution}{15.10}
若 \(E\subset D\) 且 \(a\) 是 \(E\) 的聚点，\(E\) 中每条逼近点列也是 \(D\) 中的逼近点列，故限制极限保持。反向取
\[
 f(x)=\frac{\abs x}{x},\qquad D=\R\setminus\{0\},\qquad E=(0,\infty).
\]
限制到 \(E\) 后 \(x\to0\) 的极限为 \(1\)，但在 \(D\) 上双侧极限不存在。
\end{solution}

\begin{solution}{15.11}
给定 \(\varepsilon>0\)，取局部 Cauchy 半径 \(\delta\)。由 \(x_n\to a\)，存在 \(N\) 使 \(m,n\ge N\) 时
\[
 0<\abs{x_m-a}<\delta,\qquad0<\abs{x_n-a}<\delta.
\]
于是 \(\abs{f(x_m)-f(x_n)}<\varepsilon\)，所以像数列是 Cauchy 列。
\end{solution}

\begin{solution}{15.12}
反设局部 Cauchy 条件失败。则存在 \(\varepsilon_0>0\)，使对每个 \(n\) 可选
\[
 x_n,y_n\in D,\quad
 0<\abs{x_n-a},\abs{y_n-a}<1/n,
\]
且 \(\abs{f(x_n)-f(y_n)}\ge\varepsilon_0\)。交错定义
\[
 z_{2n-1}=x_n,\qquad z_{2n}=y_n.
\]
则 \(z_n\to a\)。若 \((f(z_n))\) 是 Cauchy 列，它的相邻配对之差最终应小于 \(\varepsilon_0\)，与构造矛盾。
\end{solution}

\begin{solution}{15.13}
先用聚点性选 \(x_n\to a\)。局部 Cauchy 条件使 \(f(x_n)\) 成为实 Cauchy 列；在此处使用实数完备性，得到 \(f(x_n)\to L\in\R\)。给定 \(\varepsilon>0\)，取局部 Cauchy 条件对 \(\varepsilon/2\) 的半径 \(\delta\)，再固定充分大的 \(n_0\)，使 \(x_{n_0}\) 位于该邻域且
\[
 \abs{f(x_{n_0})-L}<\varepsilon/2.
\]
对任意同一删心邻域内的 \(x\)，
\[
 \abs{f(x)-L}
 \le\abs{f(x)-f(x_{n_0})}+\abs{f(x_{n_0})-L}
 <\varepsilon.
\]
故 \(f(x)\to L\)。
\end{solution}

\begin{solution}{15.14}
取有理数列 \(q_n\to\sqrt2\)，例如把 \(\sqrt2\) 的十进制小数截断到第 \(n\) 位，并设
\[
 D=\{0\}\cup\{1/n:n\in\Npos\},\qquad f(1/n)=q_n.
\]
因 \((q_n)\) 是有理 Cauchy 列，靠近 \(0\) 的两个定义域点对应充分靠后的两项，故局部 Cauchy 条件成立。若存在有理极限 \(r\)，则沿 \(1/n\to0\) 应有 \(q_n\to r\)，而实数中的唯一极限为 \(\sqrt2\notin\Q\)，矛盾。
\end{solution}

\begin{solution}{15.15}
若 \(f(x)\to L\) 且 \(f(x)-g(x)\to0\)，则
\[
 g(x)=f(x)-(f(x)-g(x))\to L.
\]
反向交换 \(f,g\) 即可。极限相同来自差趋于零，或直接由极限唯一性得到。
\end{solution}

\begin{solution}{15.16}
在 Cauchy 条件中取 \(\varepsilon=1\)，得到半径 \(\delta\)。由聚点性选一个固定
\[
 x_0\in D,\qquad0<\abs{x_0-a}<\delta.
\]
对任意同一删心邻域内的 \(x\)，有
\[
 \abs{f(x)}\le\abs{f(x)-f(x_0)}+\abs{f(x_0)}
 <1+\abs{f(x_0)}.
\]
故局部有界。
\end{solution}

\begin{solution}{15.17}
取 \(f(x)=\sin(1/x)\)（\(x\ne0\)）。给定 \(c\in[-1,1]\)，选 \(\theta\) 使 \(\sin\theta=c\)，并令
\[
 x_n=\frac1{\theta+2\pi n}
\]
（舍去可能使分母为零的有限项）。则 \(x_n\to0\)，且 \(f(x_n)=c\)。另一方面，取对应 \(c=1\) 与 \(c=-1\) 的两条点列，像极限不同，故函数极限不存在。
\end{solution}

\begin{solution}{15.18}
先选任意一条 \(x_n\to a\)。由条件，取 \(y_n=x_{n+1}\) 可知相邻差趋零，但这还不足以单独证明收敛。更有效的是证明 \(f\) 满足局部 Cauchy 条件。若不满足，存在 \(\varepsilon_0>0\) 及两列 \(u_n,v_n\to a\)，使
\[
 \abs{f(u_n)-f(v_n)}\ge\varepsilon_0
\]
对所有 \(n\) 成立，这与题设的差趋零矛盾。由函数 Cauchy 准则，有限极限存在。
\end{solution}

\begin{solution}{15.19}
若局部 Cauchy 条件失败，上一题构造出 \(u_n,v_n\to a\) 且配对像差不趋零。将两列交错为
\[
 z_{2n-1}=u_n,\qquad z_{2n}=v_n.
\]
则 \(z_n\to a\)。若题设对任意两列成立，可分别取交错列的奇偶子列，仍得到二者像差应趋零，与构造矛盾。于是局部 Cauchy 条件成立，再用准则。
\end{solution}

\begin{solution}{15.20}
Heine 定理不需要值域完备。其必要性只把邻域定义传给数列，充分性只从极限失败构造坏点列；在任意度量值域中都成立。函数 Cauchy 准则的必要性也不需要完备性，但充分性必须把一条像 Cauchy 列收敛到值域中的某点，因此需要完备性，或至少需要所产生的像列确实有值域内极限。
\end{solution}

\begin{solution}{15.21}
平行结构如下：
\begin{enumerate}
 \item 收敛推出 Cauchy：两种准则都用三角不等式把两项经共同极限连接。
 \item Cauchy 推出有界：固定一个尾项或固定一个删心邻域中的参考点。
 \item 构造候选极限：数列准则直接用完备性；函数准则先借聚点性选逼近点列，再对其像列使用完备性。
 \item 推广到全部对象：数列情形把某个收敛子列极限推广到全列；函数情形把参考点列的像极限推广到整个删心邻域。两步都再次用三角不等式。
\end{enumerate}
\end{solution}

\begin{solution}{15.22}
在区间定义域中，只检验从左、从右分别单调趋于 \(a\) 的点列已经足够。若极限失败，坏点数列可按距离递减抽取，使其绝对距离单调趋零；再从左右两侧之一抽取无限子列，得到单调趋近点列。

对一般稀疏定义域，仍可从任意坏点列抽取距离严格递减的子列，但点值本身未必单调；然而它必有全在左侧或全在右侧的无限子列，在单侧上距离递减就意味着点值单调。因此在实直线上，只要允许左右两个方向并要求点列严格逼近，检验全部单调逼近点列仍足够。若额外限制为某一种固定单调方向，则会漏掉另一侧。
\end{solution}

\begin{solution}{15.23}
\(\mathbf1_{\Q}\) 的振荡来自两个稠密集，选有理列与无理列即可。Thomae 型函数在无理点取零、在约分分母为 \(q\) 的有理点取 \(1/q\)；证明连续点处趋零需要控制有限个小分母有理数，证明有理点处不连续则用无理列。函数 \(\sin(1/x)\) 的振荡来自相位 \(1/x\) 无界，可精确选相位恒落在正峰与负峰的两列。三者依次体现稠密分割、分母分层和高频振荡。
\end{solution}

\begin{solution}{15.24}
设 \(Y\) 序列完备。由局部 Cauchy 条件选择 \(x_n\to a\)，像列 \(f(x_n)\) 是 \(Y\) 中的 Cauchy 列，故按序列完备性收敛到某个 \(L\in Y\)。随后用
\[
 d(f(x),L)\le d(f(x),f(x_{n_0}))+d(f(x_{n_0}),L)
\]
把参考列极限推广到整个删心邻域。证明只涉及序列，因此序列完备足够；在度量空间中，完备与序列完备等价。
\end{solution}

\section*{C组　单侧、无穷与扩充实极限}
\addcontentsline{toc}{section}{C组　单侧、无穷与扩充实极限}

\begin{solution}{16.1}
右极限 \(f(x)\to L\) 的定义是
\[
 \forall\varepsilon>0\ \exists\delta>0\ \forall x\in D,\quad
 a<x<a+\delta\Rightarrow\abs{f(x)-L}<\varepsilon.
\]
左无穷极限 \(f(x)\to+\infty\) 是
\[
 \forall M>0\ \exists\delta>0\ \forall x\in D,\quad
 a-\delta<x<a\Rightarrow f(x)>M.
\]
在 \(+\infty\) 处趋于有限 \(L\) 是
\[
 \forall\varepsilon>0\ \exists A\in\R\ \forall x\in D,\quad
 x>A\Rightarrow\abs{f(x)-L}<\varepsilon.
\]
\end{solution}

\begin{solution}{16.2}
双侧极限推出两个单侧极限，只需限制量词范围。反向分别取得对同一 \(\varepsilon\) 有效的半径 \(\delta_-,\delta_+\)，取其最小值。删心邻域中的任意点必在左侧或右侧，故相应估计成立。两侧聚点保证左右极限都不是空泛条件，也使“存在且相等”的表述有实际内容。
\end{solution}

\begin{solution}{16.3}
若 \(x\to m-\)，充分靠近时 \(m-1<x<m\)，故 \(\lfloor x\rfloor=m-1\)。若 \(x\to m+\)，充分靠近时 \(m<x<m+1\)，故 \(\lfloor x\rfloor=m\)。因此
\[
 \lim_{x\to m-}\lfloor x\rfloor=m-1,\qquad
 \lim_{x\to m+}\lfloor x\rfloor=m.
\]
\end{solution}

\begin{solution}{16.4}
当 \(x<0\) 时 \(\abs x/x=-1\)，当 \(x>0\) 时为 \(1\)。故左右极限分别为 \(-1,1\)。两者不相等，所以双侧极限不存在。
\end{solution}

\begin{solution}{16.5}
给定 \(M>0\)，取 \(\delta=1/M\)。若 \(0<x<\delta\)，则 \(x<1/M\)，从而 \(1/x>M\)。这正是右侧趋于 \(+\infty\) 的定义。
\end{solution}

\begin{solution}{16.6}
给定 \(M>0\)，要求 \(3x+1<-M\)，等价于
\[
 x<-\frac{M+1}{3}.
\]
取 \(A=-(M+1)/3\)，则所有 \(x<A\) 都满足所需不等式，故极限为 \(-\infty\)。
\end{solution}

\begin{solution}{16.7}
若定义对所有 \(M>0\) 成立，则给定 \(A\in\R\)，在 \(A>0\) 时取 \(M=A\)，在 \(A\le0\) 时取任意正阈值如 \(M=1\)，均得最终 \(f(x)>A\)。反之，“对每个实阈值”显然包含全部正阈值，故恢复原定义。
\end{solution}

\begin{solution}{16.8}
若 \(f(x)\to+\infty\)，则给定 \(M>0\)，最终有 \(f(x)>M\)，因而 \(\abs{f(x)}>M\)。反向不成立，例如 \(f(x)=1/x\) 在 \(x\to0\) 时满足 \(\abs{f(x)}\to+\infty\)，但左右符号相反，原函数没有双侧扩充实极限。
\end{solution}

\begin{solution}{16.9}
量词否定为
\[
 \exists M_0>0\ \forall\delta>0\ \exists x\in D,\quad
 0<\abs{x-a}<\delta,\quad f(x)\le M_0.
\]
按 \(\delta=1/n\) 选点，得到 \(x_n\to a\) 且 \(f(x_n)\le M_0\) 对所有 \(n\) 成立，所以像列不趋于 \(+\infty\)。反之，存在这样的见证列会与扩展 Heine 定理冲突。
\end{solution}

\begin{solution}{16.10}
给定 \(M>0\)，取 \(\delta=1/\sqrt M\)，则 \(0<\abs x<\delta\) 推出 \(1/x^2>M\)。对 \(1/x\)，沿 \(x_n=1/n\) 趋于 \(+\infty\)，沿 \(y_n=-1/n\) 趋于 \(-\infty\)，所以既无有限极限，也不趋于同一个无穷方向。
\end{solution}

\begin{solution}{16.11}
由 \(g(x)\to L\) 的局部有界性，可在某相应邻域内取 \(\abs{g(x)}\le K\)。给定 \(M>0\)，由 \(f(x)\to+\infty\) 使 \(f(x)>M+K\)。在共同邻域内
\[
 f(x)+g(x)\ge f(x)-\abs{g(x)}>M.
\]
\end{solution}

\begin{solution}{16.12}
由 \(g(x)\to M>0\)，最终有 \(g(x)>M/2\)。给定阈值 \(A>0\)，由 \(f(x)\to+\infty\) 使 \(f(x)>2A/M\)。共同邻域内
\[
 f(x)g(x)>\frac{2A}{M}\frac M2=A.
\]
\end{solution}

\begin{solution}{16.13}
在 \(x\to+\infty\) 时取
\[
 (f,g)=(x+1,x),\qquad (2x,x),\qquad
 (x+\sin x,x).
\]
每一对的两个函数都趋于 \(+\infty\)，而差依次为 \(1\)、\(x\to+\infty\)、\(\sin x\) 无极限。
\end{solution}

\begin{solution}{16.14}
在 \(x\to+\infty\) 时统一取 \(g(x)=x\)，并依次取
\[
 f_1(x)=x^{-2},\quad f_2(x)=x^{-1},\quad
 f_3(x)=x^{-1/2},\quad f_4(x)=\frac{\sin x}{x}.
\]
则 \(f_j\to0\)，而乘积依次为 \(1/x\to0\)、\(1\)、\(\sqrt x\to+\infty\)、\(\sin x\) 无极限。
\end{solution}

\begin{solution}{16.15}
必要性：给定 \(\varepsilon>0\)，由函数极限取阈值 \(A\)；由 \(x_n\to+\infty\)，充分大时 \(x_n>A\)，故像误差小于 \(\varepsilon\)。

充分性用反证。若函数极限不为 \(L\)，存在 \(\varepsilon_0>0\)，使对每个 \(A\) 都可选 \(x\in D\)、\(x>A\) 且 \(\abs{f(x)-L}\ge\varepsilon_0\)。令 \(A=n\) 逐项选择 \(x_n\)，则 \(x_n\to+\infty\)，但像列不趋于 \(L\)，矛盾。
\end{solution}

\begin{solution}{16.16}
设 \(F(x)\to L\)（\(x\to+\infty\)）。给定 \(\varepsilon>0\)，取 \(A>0\) 使 \(x>A\) 时误差小于 \(\varepsilon\)，再令 \(\delta=1/A\)。若 \(0<t<\delta\)，则 \(1/t>A\)，故 \(\abs{F(1/t)-L}<\varepsilon\)。

反之，给定 \(\varepsilon>0\)，由 \(F(1/t)\to L\)（\(t\to0+\)）取 \(\delta>0\)，令 \(A=1/\delta\)。若 \(x>A\)，则 \(0<1/x<\delta\)，故 \(\abs{F(x)-L}<\varepsilon\)。
\end{solution}

\begin{solution}{16.17}
若 \(f(x)\to0\) 且最终为正，给定 \(M>0\)，把零极限精度取为 \(1/M\)，便有
\[
 0<f(x)<1/M,\qquad1/f(x)>M.
\]
反之，若 \(1/f(x)\to+\infty\)，给定 \(\varepsilon>0\)，取阈值 \(1/\varepsilon\)，则最终
\[
 0<f(x)<\varepsilon.
\]
\end{solution}

\begin{solution}{16.18}
因 \(x^2\to0+\)（\(x\to0\)），若 \(F(u)\to L\)（\(u\to0+\)），复合极限给出 \(F(x^2)\to L\)。反向也成立：给定任意 \(u_n\to0+\)，取 \(x_n=\sqrt{u_n}\to0+\)，由 \(F(x_n^2)\to L\) 得 \(F(u_n)\to L\)，再用 Heine 定理。因而这两个极限等价；它们都不涉及 \(F\) 在 \(0\) 左侧的行为。
\end{solution}

\begin{solution}{16.19}
对有限 \(\alpha\)，邻域基可取 \((\alpha-\delta,\alpha+\delta)\)；对 \(+\infty\) 可取 \((A,+\infty]\)；对 \(-\infty\) 可取 \([-\infty,A)\)。目标 \(\beta\) 也有同样三类邻域。任取三种输入中心与三种输出中心组合，统一条件都是：每个目标邻域 \(V\) 的原像在相应意义下最终包含输入中心的某个删心邻域。九种组合正好对应有限或无穷自变量与有限或无穷函数值。
\end{solution}

\begin{solution}{16.20}
必要性对三种输入中心都由“点列最终进入每个输入邻域”推出。充分性失败时固定一个目标坏邻域 \(V\)：有限 \(\alpha\) 时在 \(0<\abs{x-\alpha}<1/n\) 中选坏点；\(+\infty\) 时在 \(x>n\) 中选坏点；\(-\infty\) 时在 \(x<-n\) 中选坏点。所得点列趋于相应 \(\alpha\)，像列却从不进入 \(V\)，与序列条件矛盾。
\end{solution}

\begin{solution}{16.21}
典型转换包括
\[
 x\to a\Longleftrightarrow h=x-a\to0,\qquad
 x\to+\infty\Longleftrightarrow t=1/x\to0+,
\]
\[
 x\to-\infty\Longleftrightarrow t=1/x\to0-,\qquad
 x\to+\infty\Longleftrightarrow t=e^{-x}\to0+.
\]
平移与倒数在相应定义域上给出等价；\(u=x^2\) 把 \(x\to0\) 映到 \(u\to0+\)，只给右极限信息；一般复合代换先给单向蕴含，只有代换对目标邻域中的允许方向具有充分覆盖性时才能反推。
\end{solution}

\begin{solution}{16.22}
例如映射
\[
 \Phi(x)=\frac2\pi\arctan x
\]
把 \(\overline{\R}\) 序同构到闭区间 \([-1,1]\)，约定 \(\Phi(\pm\infty)=\pm1\)。可定义
\[
 d_*(x,y)=\abs{\Phi(x)-\Phi(y)}.
\]
该度量的收敛恰好是扩充实数的序拓扑收敛：有限点处等价于通常收敛，趋于 \(+\infty\) 等价于 \(\Phi(x)\to1\)，负无穷类似。
\end{solution}

\begin{solution}{16.23}
令
\[
 L=\sup\{f(x):x>A\},
\]
该上确界有限。给定 \(\varepsilon>0\)，按上确界定义存在 \(x_0>A\) 使 \(f(x_0)>L-\varepsilon\)。由单调递增性，所有 \(x>x_0\) 满足
\[
 L-\varepsilon<f(x)\le L.
\]
故 \(f(x)\to L\)（\(x\to+\infty\)）。完备性通过上确界的存在进入证明。
\end{solution}

\begin{solution}{16.24}
当 \(x\to+\infty\) 时，
\[
 (x+1)-x\to1,\quad 2x-x\to+\infty,\quad
 x-(2x)\to-\infty,
\]
而 \((x+\sin x)-x=\sin x\) 无极限。对 \(0\cdot\infty\)，取无穷因子 \(g=x\)，零因子分别取 \(x^{-2},x^{-1},x^{-1/2},\sin x/x\)，乘积依次趋于 \(0,1,+\infty\) 或无极限。这些结果说明未定式没有预先指定数值。
\end{solution}
